% Dexterous robotic hands, which mimic the anthropomorphic structure of
% the human hand, typically possess a large number of degrees of freedom
% (DoFs). This high level of dexterity makes them well suited for
% complex and contact-rich tasks such as grasping (e.g., pinching,
% tripod, and power grasps) and in-hand manipulation (e.g., reorientation
% and twisting). However, this same dexterity also introduces mechanical
% delicacy and control complexity, making dexterous hands among the most
% fragile components in robotic systems, particularly in highly dynamic
% platforms such as humanoid robots. 

Dexterous robotic hands typically possess a large number of degrees of freedom (DoFs) 
and exhibit complex contact dynamics, enabling them to perform fine-grained and diverse 
manipulation behaviors, such as multi-finger grasping~\cite{brahmbhatt2019contactgrasp}, in-hand 
reorientation~\cite{chen2022system,andrychowicz2020learning}, and precise control~\cite{luo2025precise}. However, 
the high-dimensional action space and intricate interaction dynamics substantially increase 
the difficulty of control and policy learning. In recent years, reinforcement learning (RL)-based 
approaches~\cite{rajeswaran2017learning,andrychowicz2020learning,akkaya2019solving,petrenko2023dexpbt} 
have achieved remarkable progress in dexterous grasping and manipulation tasks. Leveraging large-scale 
parallel simulation and high-fidelity physics engines, RL enables the automatic learning of complex
control policies in simulated environments without the need for manually designed controllers. 
These methods have demonstrated strong generalization capabilities across multi-object and 
multi-task scenarios. 

Despite these advances, existing reinforcement learning frameworks primarily optimize for task 
success rates and rarely incorporate explicit modeling or constraints related to safety-critical 
behaviors. During training, common failure modes include severe collisions between the hand and 
the environment, unintended object ejection due to excessive contact forces, and loss of stable 
hand posture~\cite{liu2023safe,yu2022dexterous}. Such unsafe behaviors not only reduce training 
efficiency but also introduce potential risks during real-world deployment. Therefore, how to 
reduce unsafe exploration while maintaining learning efficiency remains a critical challenge in 
dexterous manipulation.

% Among the central challenges in dexterous manipulation research is how 
% to acquire reliable skills for complex systems like dexterous hands and 
% how to bridge the sim to real gap effectively. 

% Recent research has made significant progress in acquiring various skills
% through learning-based approaches. Policies can be learned
% via reinforcement learning in simulated environments conditioned with task
% specifics. Concurrent advances in vectorized simulation and increasingly realistic
% physics engines equipped with a wide range of accurate sensor choices,
% have greatly simplified the reinforcement learning process.

%====================================================
% To mitigate the performance gap between simulation and real-world deployment, policy distillation has become a widely adopted paradigm in recent years. The core idea is to first train a high-performing expert policy in simulation using privileged state information, and then distill its decision-making behavior into a student policy that relies solely on visual observations or partially observable inputs. This process enables modality transfer from privileged state information to perception-based inputs, facilitating efficient real-world deployment. 
% Policy distillation reduces reliance on accurate state estimation and allows knowledge integration from multiple teacher policies, thereby significantly improving execution stability and generalization performance in real-world environments. Consequently, it has gained increasing attention in dexterous grasping and complex manipulation tasks. 
% However, existing distillation methods~\cite{} typically rely on behavior cloning or the standard DAgger framework for data aggregation. Their primary objective is to mitigate performance degradation caused by distribution shift, rather than explicitly enforcing safety constraints during student policy learning. During distillation, the student policy may deviate from the teacher’s state–action distribution and enter regions that are either insufficiently covered by the teacher or inherently high-risk, leading to unsafe behaviors such as severe collisions or unstable hand configurations. 
% Therefore, introducing explicit safety mechanisms into the distillation process—so that high-risk exploration can be avoided without sacrificing learning efficiency—remains an important and largely underexplored problem. Existing safety-aware distillation approaches~\cite{}, such as SafeDAgger~\cite{}, have primarily focused on autonomous driving scenarios and have not been systematically studied in the context of dexterous manipulation.

To mitigate the performance gap between simulation and real-world deployment,
policy distillation has become a widely adopted paradigm in recent years. The
core idea is to first train a high-performing expert policy in simulation using
privileged state information, and then distill its decision-making behavior into
a student policy that relies solely on visual observations or partially
observable inputs. This process enables modality transfer from privileged state
information to perception-based inputs, facilitating efficient real-world
deployment.
Policy distillation reduces reliance on accurate state estimation and allows
knowledge integration from multiple teacher policies, thereby substantially
improving execution stability and generalization performance in real-world
environments. Consequently, it has gained increasing attention in dexterous
grasping and complex manipulation tasks.
However, existing distillation methods~\cite{ross2011reduction,singh2024dextrah}
typically rely on behavior cloning or the standard DAgger framework for data
aggregation. Their primary objective is to mitigate performance degradation
caused by distribution shift, rather than explicitly enforcing safety
constraints during student policy learning. During distillation, the student
policy may deviate from the teacher’s state--action distribution and enter
regions that are either insufficiently covered by the teacher or inherently
high-risk, leading to unsafe behaviors such as severe collisions or unstable
hand configurations.
Therefore, introducing explicit safety mechanisms into the distillation
process---so that high-risk exploration can be avoided without sacrificing
learning efficiency---remains an important and largely underexplored problem.
Existing safety-aware distillation
approaches, such as
SafeDAgger~\cite{zhang2016query}, have primarily focused on autonomous driving
scenarios and have not been systematically studied in the context of dexterous
manipulation.

%===================================================================================
% To address the sim-to-real transfer problem, techniques such as domain
% randomization and domain adaptation have been widely adopted. However,
% these methods often lack interpretability. In contrast, imitation learning
% has played a key role in multi-policy compression and efficient knowledge
% transfer. Among these approaches, policy distillation has attracted
% increasing research attention not only for its ability to consolidate
% multiple specialized teacher skills, but also for its potential to enable
% observation modality transfer.

% Most prior work studies policy distillation purely in simulation, where it is
% used for behavior replication or policy compression. However, both in simulation
% and in real-world deployment, student policies are limited by imperfect
% imitation and are inevitably exposed to distributional shift, visiting states
% that were not encountered during the teacher's learning process. Under such
% conditions, safe exploration of the task space is critical to avoid catastrophic
% outcomes.

% To address this challenge, we propose a pipeline that combines
% reinforcement-learning-based teacher training for per-object grasping skills
% with distillation into an RGB-based student policy. We further enhance the
% distillation process with a teacher-intervened exploration strategy, guided by a
% combined decision signal derived from both teacher--student disagreement and a
% risk-prediction probability.

The contributions of this work are summarized as follows:
\begin{itemize}
    \item We propose a safety-sensitive framework for distilling RL-acquired teacher policies into a single RGB-based
    student policy. To the best of our knowledge, this is the first attempt in dexterous manipulation to incorporate
    a mixed intervention mechanism that combines teacher--student action
    disagreement with a Bellman success-value critic trained from lift-success
    labels. Low predicted success triggers proactive teacher takeover.

    \item We propose a safety-aware RL reward design that explicitly accounts for risky events in 
    arm--hand control scenarios, including harmful collisions, object escape, palm flipping, and the hand moving 
    too far from the workspace, thereby encouraging safer exploration during teacher training.
    
    \item We empirically demonstrate that combining disagreement and success-value prediction substantially reduces unsafe
    exploration compared with vanilla DAgger, achieving a relative reduction of \textbf{80\%} in unsafe exploration
    rate in the early training phase. We also observe a \textbf{5\%} improvement in unsafe exploration rate compared
    with disagreement-only SafeDAgger.
\end{itemize}
